\documentclass[12pt,oneside]{book}
\usepackage[brazil]{babel}
\usepackage[utf8]{inputenc}
\usepackage[T1]{fontenc}

%\usepackage{DejaVuSerif}		% Fonte do texto

%%%%%%%%%%%%%%%%%%%%% PACOTES USADOS PARA FAZER A CAPA
%

\usepackage[a4paper,bottom=2cm]{geometry} % geometry/margins of page
\savegeometry{origin}
\geometry{rmargin=2cm,lmargin=2cm}% for the title page
\usepackage{multicol}		% multicolumn environments
\usepackage{tikz}			% used for the 'logo' (remove if unwanted)
\usepackage{subfig}			% Usada para alinhar o logo da capa
%
%%%%%%%%%%%%%%%%%%%%%%


%\usepackage{imakeidx}\makeindex
\usepackage{amsmath,amsthm,amssymb,color}
\usepackage[Lenny]{fncychap} % Usado para decorar Headers and footers
%
%
\usepackage{mathrsfs}
\usepackage{hyperref}


% Muda de Chapter ou Capítulos para Aulas
\addto\captionsbrazil{\renewcommand{\chaptername}{Aula}}




%%%%%%%%%%%%% INICIO - Alguns Comandos Úteis
\newcommand{\C}{\mathbb{C}}
\newcommand{\R}{\mathbb{R}}
\newcommand{\Q}{\mathbb{Q}}
\newcommand{\Z}{\mathbb{Z}}
\newcommand{\N}{\mathbb{N}}
\renewcommand{\P}{\mathbb{P}}
\newcommand{\F}{\mathcal{F}}



\newtheorem{teorema}{Teorema}[chapter]
\newtheorem{proposicao}[teorema]{Proposição}
\newtheorem{lema}[teorema]{Lema}
\newtheorem{corolario}[teorema]{Corolário}
\newtheorem{definicao}[teorema]{Definição}
\newtheorem{observacao}[teorema]{Observação}
\newtheorem{exemplo}[teorema]{Exemplo}
\newtheorem{exercicio}[teorema]{Exercício}
%
\renewenvironment{proof}[1][Demonstração]
{\noindent\textbf{#1. \ }}{\hspace{\stretch{1}}\rule{1ex}{1ex}}
%


% operadores matemáticos
\DeclareMathOperator{\e}{e}
\DeclareMathOperator{\vol}{vol}
\DeclareMathOperator{\Id}{Id}
\DeclareMathOperator{\Graf}{Graf}
\DeclareMathOperator{\intt}{int}
\DeclareMathOperator{\sen}{sen}
\renewcommand{\ker}{Nuc}
\DeclareMathOperator{\supp}{supp}

%%%%%%%%%%%%% FIM    - Alguns Comandos Úteis



%   Comandos para gerar texto em cores
\definecolor{Red}{cmyk}{0,1,1,0}
\def\red{\color{Red}}
\definecolor{Blue}{cmyk}{1,1,0,0}
\def\blue{\color{Blue}}
%  



% Novos Comandos
\newcommand{\tn}{\textnormal}



% Begin document
% ---------------
\begin{document}

%%%%%%%%%%%%%%%%%%  Capa
\frontmatter 
%%%% Este arquivo não deve ser compilado
%%%% Ele possui as informações necessárias 
%%%% para fazer a capa das notas de aulas
%
%
%
\begin{titlepage}			% environment specially for titlepages

\begin{center}				% center everything on the titlepage

	% title	
	{\fontsize{16mm}{11mm}
		\selectfont
		\textbf{Probabilidade} 
		\\[0.3cm]
		\textbf{e}
		\\[0.3cm] 
		\textbf{Teoria da Medida} 
		}

	\vspace{60mm}				% vertical spacing


	% subtitle 1
	{\fontsize{14pt}{14pt}\selectfont
		\textbf{Notas de Aula}
		}
	
	\vspace{7pt}				% vertical spacing	


	% subtitle 2	
	{\fontsize{14pt}{14pt}\selectfont
		\textbf{Julho / 2014}
		}
	
	\vfill						% vertical spacing	





%%%%%%%%%%%%%%%%%%%% BEGIN Authors
	\begin{multicols}{3}		% start two column layout

	% author 1 (left)
	{\fontsize{12pt}{12pt}\selectfont
		\textbf{\MakeUppercase{Josimar Aguirre}}
		}
		MAT/UnB
	\columnbreak				% switch over to second column


	% author 2 (midle)
	{\fontsize{12pt}{12pt}\selectfont
		\textbf{\MakeUppercase{Eduardo Antonio}}
		}
		MAT/UnB
	\columnbreak				% switch over to second column


	% author 3 (left)
	{\fontsize{12pt}{12pt}\selectfont
		\textbf{\MakeUppercase{Leandro Cioletti}}
		}
		MAT/UnB	
	\end{multicols}				% end two column layout
%%%%%%%%%%%%%%%%%%%%  END Authors





%%%%%%%%%%%%%%%%%%%%  Logomarca
	\vspace{10mm}				% vertical spacing	
	\begin{figure}[h]            % Logo da UnB
		\centering
			\subfloat{\includegraphics[scale=1.0]{Figuras/logo-unb.pdf}}
			\subfloat{\raisebox{0.15cm}{\fontsize{25pt}{20pt}\selectfont \ UnB}}
	\end{figure} 

\end{center}
\end{titlepage}
\loadgeometry{origin}% restore the orign margin setting



%%%%%%%%%%%%%%%%%%%%  Índice
\tableofcontents


 
%%%%%%%%%%%%%%%%%%%% Conteúdo das aulas    
\mainmatter 

% Aula 1 
\chapter[Aula 1]{Conjuntos e Limite de Sequências de Conjuntos}
\chaptermark{}

Nesta seção são introduzidos alguns fatos básicos da teoria
conjuntos. O ponto de vista adotado aqui é o da "teoria ingênua
de conjuntos" e não temos objetivos de discutí-la do ponto 
vista axiomático.

  
\section{Teoria Básica de Conjuntos}
\textbf{Notações}.
%
Em geral, vamos trabalhar em um espaço que será denotado 
por $\Omega$. Assim, operações de conjuntos serão sempre 
consideradas com relação a este espaço. 
A coleção de todos os subconjuntos de $\Omega$ será 
denotada por $\mathcal{P}(\Omega)\equiv \{A: A\subset \Omega\}$ e 
será chamado de conjunto das partes de $\Omega$.
Em geral, usaremos as letras maiúsculas $A,B$ e etc. para denotar
um subconjunto arbitrário de $\Omega$. Quando quisermos nos referir
a uma coleção de subconjuntos de $\Omega$ usaremos letras maiúsculas 
caligráficas como $\mathcal{A}, \mathcal{B}$ e etc. Finalmente,
usaremos, na maioria das vezes, a notação $\omega\in\Omega$ 
para denotar um ponto do espaço $\Omega$.


\noindent
\textbf{Operações de Conjuntos}

O complementar de um conjunto $A$ será 
denotado por $A^{c}= \{w \in \Omega; w \notin A\}$. 
Se $I$ é um conjunto de índices arbitrários e 
para cada $i\in I$ temos que $A_i\subset \Omega$ então 
a interseção e união dos conjuntos $A_i$'s sobre a coleção
$I$ são dados, respectivamente, por 
\begin{enumerate}
	\item[]  $\displaystyle\bigcap_{t \in T}{A_t}
							= 
						\{ w \in \Omega;\ w \in A_t,\ \forall t \in T\}$.
	\item[]  $\displaystyle\bigcup_{t \in T}{A_t}
							= 
						\{ w \in \Omega;\ w \in A_t,\ \text{para algum}\ t \in T \}$.
\end{enumerate}
%
%
%
Se $A \cap B = \emptyset$, dizemos que $A$ e $B$ são disjuntos.
Dizemos que uma sequência de conjuntos 
$A_1, A_2,\ldots$ é mutuamente disjunta (ou dois a dois disjunta) 
se $A_i \cap A_j = \emptyset$ sempre que $i \neq j$. 
Quando $A_1,A_2,\ldots $ for uma sequência mutuamente disjunta, 
vamos usar a notação abaixo para denotar a união dos conjuntos $A_n$'s
$$
	\displaystyle\bigcup_{n \geqslant 1}{A_n} = \sum_{n \geqslant 1} A_n.	
$$

\noindent
\textbf{Outras Notações:} Quando for conveniente 
usaremos as seguintes notações e convenções:
\begin{enumerate}
\item[$\blacklozenge$] $AB\equiv A \cap B$.
\item[$\blacklozenge$] $A \setminus B \equiv A \cap B^c$.
\item[$\blacklozenge$] Diferença Simétrica $A\triangle B = (A \setminus B) \cup (B\setminus A).$
\item[$\blacklozenge$] $\emptyset^c = \Omega$.
\item[$\blacklozenge$] $\Omega^c = \emptyset$.
\end{enumerate}




\begin{exercicio}[Associatividade] Mostre que para quaisquer subconjuntos $A,B$ e $C$ 
de $\Omega$ valem as seguintes igualdades:
\begin{enumerate}
	\item $(A\cup B)\cup C= A \cup ( B \cup C)$.
	\item $(A\cap B)\cap C= A \cap ( B \cap C)$.
\end{enumerate}
\end{exercicio}






\begin{exercicio}[Leis de de Morgan] Sejam $I$ um conjunto arbitrário de índices
e $A_i\subset \Omega$ para todo $i\in I$. Mostre que
as seguintes igualdades são válidas:
\begin{enumerate}
\item 
$
\left( \displaystyle\bigcup_{i \in I}{A_i} \right)^c 
= 
\displaystyle\bigcap_{i \in I}{{A_i}^c}
$.

\item
$
\left( \displaystyle\bigcap_{i \in I}{A_i} \right)^c 
= 
\displaystyle\bigcup_{i \in I}{A_i^c}
$.
%
\end{enumerate}
\end{exercicio}





\begin{exercicio}[Distributiva] Sejam $I$ um conjunto arbitrário de índices, 
$B\subset \Omega$ e $A_i\subset \Omega$ para todo $i\in I$. 
Mostre que as seguintes igualdades são válidas:
%
\begin{enumerate}
\item 
$
B \cap \left( \displaystyle\bigcup_{i \in I}{A_i} \right) 
= 
\displaystyle\bigcup_{i \in I}{(B\cap A_i)} 
$.
%
\item
$
B \cup \left( \displaystyle\bigcap_{i \in I}{A_i} \right) 
= 
\displaystyle\bigcap_{i \in I}{(B\cup A_i)} 
$.
\end{enumerate}
%
\end{exercicio}







\begin{definicao}[Função Indicadora]\label{def-funcao-indicadora}
Seja $A \subseteq \Omega$. 
A função indicadora \index{função!indicadora} 
de $A$ é denotada por $1_A: \Omega \to \R$ e definida 
como segue
$$
	1_A(w) =
	\begin{cases}
		1, & \text{se}\ w \in A; \\
		0, & \text{caso contrário.}
	\end{cases}
$$
\end{definicao}








\begin{observacao} Da definição de função indicadora, podemos mostrar facilmente 
as seguintes relações,
as quais serão importantes ao longo do texto
\begin{enumerate}
\item 
$1_A \leqslant 1_B \Leftrightarrow A \subseteq B$.

\item
$1_{A^c}= 1- 1_A$.
\end{enumerate}
\end{observacao}

\vspace*{1cm}
\textbf{Limite de Conjuntos}

Seja uma sequência de subconjuntos $A_n \in \Omega$, 
vamos a definir as noções de limite para estos conjuntos como segue
\begin{enumerate}
\item[$\blacklozenge$] $\inf \limits_{k\geqslant n} A_k = \displaystyle\bigcap_{k=n}^{\infty}{A_k}$.
\item[$\blacklozenge$] $\sup \limits_{k\geqslant n} A_k = \displaystyle\bigcup_{k=n}^{\infty}{A_k}$.
\item[$\blacklozenge$] 
	$\liminf \limits_{n \to \infty} A_n 
	= 
	\displaystyle\bigcup_ {n\geqslant 1} \left(\displaystyle\bigcap_{k=n}^{\infty}{A_k} \right)$.
\item[$\blacklozenge$] 
	$\limsup \limits_{n \to \infty} A_n 
	= 
	\displaystyle\bigcap_ {n\geqslant 1} \left(\displaystyle\bigcup_{k=n}^{\infty}{A_k} \right)$.
\end{enumerate}

\begin{definicao}
	Se para alguma sequência de conjuntos $B_n \subseteq \Omega$ temos

	$$
		\liminf \limits_{n \to \infty} B_n = \limsup \limits_{n \to \infty} B_n = B,
	$$	
	dizemos que existe $\lim \limits_{n \to \infty} B_n =B$.
\end{definicao}

Note que $\liminf \limits_{n\to \infty}A_n 
			=
			\lim \limits_{n\to \infty}\left(\inf \limits_{k\geqslant n}A_k \right)
		 $.


\begin{lema}
Seja $\{A_n\}$ uma sequência de subconjuntos de $\Omega$.
\begin{enumerate}
\item[a)] $\limsup A_n = \{ w \in \Omega;\ \sum_{n \geqslant 1} 1_{A_n}(w)= \infty \}$.
\item[b)] 
\begin{align*}
\vspace*{-2cm}
\liminf A_n & = \{ w \in \Omega;\ w \in A_n \text{ para todo $n$ excepto uma quantidade finita} \} \\
			& = \{w \in \Omega;\ \sum_{n\geqslant 1} 1_{A_n^c}(w) < \infty \} \\
			& = \{w \in \Omega;\ w \in A_n, \forall n \geqslant n_0(w) \}. 
\end{align*} 
\end{enumerate}
\end{lema}

\begin{proof}
Vejamos o item $(a)$. Suponha que $w \in \limsup A_n$, então $w \in \displaystyle\bigcup_{k \geq n} A_k,
\forall n \in \N$, logo existe $k_n \geqslant n$ tal que $w \in A_{k_n}$. Assim

	$$
		\sum \limits_{k \geqslant 1} 1_{A_k}(w) \geqslant \sum \limits_{n\geqslant 1} 1_{A_{k_n}}(w) 
		= 
		\infty.
	$$

Reciprocamente, se $ w \in \{ w \in \Omega;\ \sum_{n \geqslant 1} 1_{A_n}(w)= \infty \}$ daí existem 
infinitos $k$ tais que $w \in A_k$. Portanto $w \in \limsup A_n$.

O item $(b)$ do lema segue direto do item $(a)$. 
\end{proof}

\begin{observacao}

\begin{enumerate}
\item[i)] $\liminf A_n \subseteq \limsup A_n$.
\item[ii)] $\left( \liminf A_n \right)^c = \limsup A_n^c$. 
\item[iii)] $\{X_n,\ n\geqslant 0\}$ sequência de v.a. No caso que desejamos mostrar que $X_n \to X$ q.c.
podemos usar o seguinte critério:

$$
	\mathbb{P}( |X_n-X| > \epsilon \ \text{infinitas vezes} )=0
$$ 
isto é, fazendo $A_n= \{ |X_n-X|> \epsilon\}$ então $\mathbb{P} (\limsup A_n )=0$.
\end{enumerate}
\end{observacao}

\vspace*{1cm}
\textbf{Sequências Monótonas}

Se $\{A_n\}$ é uma sequência de conjuntos de $\Omega$. Dizemos que $\{A_n\}$ é monótona não
decrescente se $A_1 \subseteq A_2 \subseteq \ldots$ Analogamente definimos sequência não decrescente.

Notação: $A_n \nearrow$ ou $A_n \uparrow$ (analogamente $A_n\searrow$ ou $A_n \downarrow$).

\begin{proposicao}
 Suponha que $\{A_n\}$ é uma sequência monótona.
 \begin{enumerate}
 \item[1)] Se $A_n \nearrow$ então 
 	$\exists \lim \limits_{n \to \infty} A_n= \displaystyle\bigcup_{n\geqslant 1} {A_n}$.
 \item[2)] Se $A_n \searrow$ então 
 	$\exists \lim \limits_{n \to \infty} A_n= \displaystyle\bigcap_{n\geqslant 1} {A_n}$.
 \end{enumerate}
\end{proposicao}


\begin{proof}
Vejamos a primeira parte. Queremos mostrar que 
$\liminf A_n = \limsup A_n =\displaystyle\bigcup_{n\geqslant 1} {A_n}$. Já que $A_j \subseteq A_{j+1}$
então $\displaystyle\bigcap_{k\geqslant n} {A_k}=A_n$.

Daí $\liminf A_n = \displaystyle\bigcup_{n\geqslant 1} {A_n}$. Agora 

$$
	\limsup A_n 
	=
	\displaystyle\bigcap_{n\geqslant 1} \left(\displaystyle\bigcup_{k\geqslant n}{A_k} \right)
	\subseteq 
	\displaystyle\bigcup_{k\geqslant 1} {A_k} 
	=
	\liminf A_n \subseteq \limsup A_n.
$$
Para a segunda parte é somente usar as leis de De Morgan. Isto conclui a prova.
\end{proof}

\vspace*{1cm}
\textbf{Relações de ``Dualidade"}
\begin{enumerate}
\item[1)] $1_{\inf_{k\geqslant n} A_k} = \inf \limits_{k\geqslant n} 1_{A_k}$.
\item[2)] $1_{\sup_{k\geqslant n} A_k} = \sup \limits_{k\geqslant n} 1_{A_k}$.
\item[3)] $1_{\cup_{n\geqslant 1} A_n} \leqslant \sum \limits_{n \geqslant 1} 1_{A_n}$.
\item[4)] $1_{\limsup A_n} = \limsup 1_{A_n}$.
\item[5)] $1_{\liminf A_n} = \liminf 1_{A_n}$.
\item[6)] $ 1_{A \triangle B} = 1_A + 1_B (\text{mod 2})$.
\end{enumerate}


% Aula 2
\chapter[Aula 2]{Espaços de Medida e o Teorema da Extensão de Caratheodóry}
\chaptermark{}
\section{Espaço de Medidas}


Seja $\Omega$ um conjunto não vazio.
Uma coleção $\mathcal{A}$ de subconjuntos de $\Omega$ é chamada de
\index{Álgebra! de conjuntos} {\it álgebra} se satisfaz as seguintes condições:
\begin{itemize}
	\item[1)] $\emptyset\in\mathcal{A}$ e $\Omega\in\mathcal{A}$;
	\item[2)] $A\in \mathcal{A} \Rightarrow A^c\in\mathcal{A}$;
	\item[3)] $A,B\in \mathcal{A} \Rightarrow A\cup B\in\mathcal{A}$.
\end{itemize}
Note que 1) e 2) implicam que $\mathcal{A}$ é fechada para uniões 
e interseções finitas. Se substituímos 3) por:
\begin{itemize}
	\item[3')] $A_n\in\mathcal{A}\ \, (n=1,2,\ldots) 
				\Rightarrow 
				\displaystyle\bigcup_{n=1}^{\infty} A_n \in \mathcal{A}$.
\end{itemize}
dizemos que $\mathcal{A}$ é uma $\sigma$-{\it álgebra} \index{$\sigma$-álgebra} . Observe que a condição 3') 
implica na condição 3) e também que uma $\sigma$-álgebra é fechada para
interseções enumeráveis.

\begin{definicao}[Medida]
Uma função $\mu:\mathcal{A}\to [0,\infty]$ é chamada de uma {\it medida sobre uma álgebra} 
$\mathcal{A}$ se 
\begin{itemize}
	\item[1)] $\mu(\emptyset)=0$;
	\item[2)] para toda sequência de conjuntos $A_n \ \ (n=1,2,\ldots)$ dois a dois disjunta
				tal que se $\cup_{n=1}^{\infty}A_n \in \mathcal{A}$, então temos
			 $\mu(\cup_{n=1}^{\infty} A_n ) =\sum_{i=1}^{\infty}\mu(A_n)$.
\end{itemize}
\end{definicao}
A segunda propriedade é conhecida como $\sigma$-aditividade e ela implica 
em aditividade finita (tomando $A_n=\emptyset$ para $n\geq m$, para algum $m\in\mathbb{N}$).


\begin{exercicio}
 Mostre que se $\Omega=\mathbb{N}$ e $\mathcal{A}= \mathcal{P}(\Omega) $ é a $\sigma$-álgebra das partes de 
 $\Omega$, então a função $\mu:\mathcal{A}\to [0,+\infty]$ cuja a imagem de um subconjunto 
 $E\subset\Omega$ é dada por 
 \[
     \mu(E)=\sharp E
 \]
 é uma medida em $\mathcal{A}$. Esta medida é chamada de medida da contagem em $\mathbb{N}$.
\end{exercicio}

\begin{exercicio}
	Seja $\Omega$ um conjunto enumerável não vazio, $\mathcal{A}=\mathcal{P}(\Omega)$ e
	 $f:\Omega\to [0,\infty]$ 
	uma função arbitrária. Mostre que a fórmula 
	\[
	    \mu(E)=\sum_{x\in E} f(x),
	\]
	determina uma medida em $\mathcal{A}$.
\end{exercicio}


\begin{definicao}[Medida de Probabilidade]
Seja $\mathcal{F}$ uma $\sigma$-álgebra de conjuntos de um espaço $\Omega$ e 
$P:\mathcal{F}\to [0,\infty]$ uma medida tal que $P(\Omega)=1$.
Neste caso dizemos que $P$ é uma medida de probabilidade. 
Normalmente o espaço $\Omega$
é chamado de espaço amostral, um elemento $E\in\mathcal{F}$ é chamado
de evento e $P(E)$ é a probabilidade de ocorrer $E$.
\end{definicao}


Sejam $\mathcal{A}$ uma álgebra e $\mu:\mathcal{A}\to [0,\infty]$ uma medida.
Se $A_n\in\mathcal{A}$ ($n=1,2,\ldots$) e
$\cup_{n=1}^{\infty} A_n\in\mathcal{A}$ 
então para todo $A\subset \cup_{n=1}^{\infty} A_n$ tal que $A\in\mathcal{A}$, temos 
$$
\mu(A) \leq \sum_{i=1}^{\infty} \mu(A_n).
$$
Para verificar que este fato é verdadeiro, definimos $B_1=A_1$
e $B_n = A_1^c\cap \ldots \cap A_{n-1}^c\cap A_n$ para todo $n\in\mathbb{N}$.
então $B_n$ forma uma sequência de conjuntos dois a dois disjuntos e além do
mais $\cup_{n=1}^{\infty} A_n=\cup_{n=1}^{\infty} B_n$, logo 
$$
\mu(A) 	= \mu(A\cap [\cup_{n=1}^{\infty} B_n])
		= \sum_{n=1}^{\infty} \mu(A\cap B_n)
		\leq \sum_{n=1}^{\infty} \mu(A_n),
$$
na última desigualdade usamos que $B_n\subset A_n$ para todo $n\in\mathbb{N}$
e monotonicidade de $\mu$. 


\begin{proposicao}[Continuidade da Medida]
Seja $\mathcal{F}$ uma $\sigma$-álgebra de conjuntos de um espaço $\Omega$ e 
$\mu:\mathcal{F}\to [0,\infty]$ uma medida. 
\begin{itemize}
\item Para qualquer sequência crescente $A_n$
($n=1,2,\ldots$) em $\mathcal{F}$, isto é,  
$A_1\subset A_2\subset\ldots$ temos que 
$$
\mu\left( \bigcup_{i=1}^{\infty} A_n \right) = \lim_{n\to\infty} \mu(A_n).
$$ 
\item Se $A_n$ é uma sequência decrescente, isto é, 
$A_1\supset A_2\supset\ldots$ e $\mu(A_1)<\infty$, então 
$$
\mu\left( \bigcap_{i=1}^{\infty} A_n \right) = \lim_{n\to\infty} \mu(A_n).
$$ 
\end{itemize}
\end{proposicao}
\begin{proof}
A prova desta proposição é deixada como exercício para o leitor.
\end{proof}


\medskip
Seja $\mu$ uma medida em uma $\sigma$-álgebra $\mathcal{F}$.
A $\sigma$-álgebra $\mathcal{F}$ é dita $\mu$-{\bf completa} 
\index{$\sigma$-álgebra!$\mu$-completa}
se para todo $N\in \mathcal{F}$ e para todo $A\subset N$, 
com $\mu(N)=0$ temos $A\in \mathcal{F}$.
Em outras palavras, $\mathcal{F}$ é $\mu$-completa 
se os subconjuntos dos conjuntos de medida $\mu$ zero
pertecem a $\sigma$-álgebra $\mathcal{F}$.
Neste caso dizemos também que $\mu$ é completa.
Dada qualquer medida $\mu$ definida sobre uma 
$\sigma$-álgebra $\mathcal{F}$ é fácil verificar que 
a coleção 
$$
\overline{\mathcal{F}}
=
\{ 	
	C=A\cup B: A\in\mathcal{F}, B\subset N\ \text{para algum}\ N\in\mathcal{F}\ \text{com}\ \mu(N)=0	
\}
$$
é uma $\sigma$-álgebra $\overline{\mu}$ completa, onde $\overline{\mu}(A\cup B):=\mu(A)$.
Observe $\overline{\mu}$ está bem definida e 
que $\overline{\mu}$ é uma medida definida em $\overline{\mathcal{F}}$ que estende $\mu$.
Esta extensão de $\mu$ é chamada de \index{Completamento! de uma medida} 
{\it completamento}
de $\mu$.





Passamos agora para o estudo de um dos teoremas básicos mais importantes 
na Teoria da Medida. 
Uma das conclusões mais importantes deste teorema, conhecido na literatura como 
Teorema de Extensão de Carathéodory, 
é que qualquer medida $\mu$ definida sobre uma álgebra
$\mathcal{A}$ sempre pode ser estendida 
a uma medida $\overline{\mu}$, 
definida sobre a
$\sigma$-álgebra gerada por $\mathcal{A}$, que 
é simplesmente a ``interseção'' de todas as $\sigma$-álgebras 
contendo $\mathcal{A}$, veja os dois exercícios abaixo.

\begin{exercicio}
	Sejam $I$ um conjunto arbitrário de índices e 
	$\mathcal{F}_i$, para cada $i\in I$,
	uma $\sigma$-álgebra de conjuntos de $\Omega$.
	Defina 
	$$
	\bigcap_{i\in I}\mathcal{F}_i := 
	\{F\subset \Omega: F\in\mathcal{F}_i, \ \forall i\in I\}.
	$$
	Mostre que $\cap_{i\in I}\mathcal{F}_i$ é uma $\sigma$-álgebra.
\end{exercicio}


\begin{exercicio}\label{exercicio-sigma-alg-gerada}
	Seja $\mathcal{A}$ uma coleção arbitrária de subconjuntos de $\Omega$.
	Mostre que 
	$$
	\bigcap_{\substack{ \mathcal{F}\supset \mathcal{A}\\[0.1cm] \mathcal{F}\ \text{é}\ 
	\sigma\text{-álgebra}}}
	 \!\!\!\!\!\!\!\!\! \mathcal{F}
	$$
	é uma $\sigma$-álgebra. Esta $\sigma$-álgebra é conhecida como a 
	$\sigma$-álgebra
	gerada pela coleção $\mathcal{A}$. Em um certo sentido, esta é a ``menor'' 
	$\sigma$-álgebra contendo a coleção $\mathcal{A}$.
\end{exercicio}



\section{O Teorema de Dynkin}

\begin{definicao} 
    Dada uma coleção $\mathcal{C} $ de subconjuntos de $\Omega$, a menor
    $\sigma$-álgebra que contém todos os eventos de $\mathcal{C}$ é chamada
    de \textbf{$\sigma$-álgebra gerada por $\mathcal{C}$} e será
    denotada por $\sigma(\mathcal{C})$.
\end{definicao} 
A existência desta $\sigma$-álgebra é garantida pela afirmação 
contida no Exercício \ref{exercicio-sigma-alg-gerada}.
	    
	    
	    
	    
	    
    \begin{exercicio}\label{exer-sigma-gerada-contida-todas}
        Seja $\mathcal{C}$ uma coleção de subconjuntos de $\Omega$ e $\mathcal{F}$
        uma $\sigma$-álgebra de $\Omega$ tal que $\mathcal{C} \subset \mathcal{F}$.
        Mostre que  $\sigma(\mathcal{C}) \subset \mathcal{F}$.
    \end{exercicio}

    
    Observamos que o exercício acima não vale para a união de $\sigma$-álgebras. 
    Dada uma coleção de $\sigma$-álgebras, digamos $\mathcal{F}_{i}, i \in I$,
    consideramos a $\sigma$-álgebra mais grossa contendo todas as $\sigma$-álgebras 
    $\mathcal{F}_{i}$, isto é, 
    $\vee_{\lambda \in \Lambda} \mathcal{F}_{\lambda} = 
    \sigma (\cup_{\lambda \in \Lambda} \mathcal{F}_\lambda)$. 
    
    É muito comum trabalhar com $\sigma$-álgebras 
    $\mathcal{F}= \sigma(\mathcal{C})$ em que $\mathcal{C}$ é uma 
    coleção fechada para a interseção finita. Esta observação motiva 
    a seguinte definição:
    
\begin{definicao}[$\pi$-sistema]\label{def-pi-sistema}
   Uma coleção $\mathcal{C} $ fechada para interseções finitas é 
    chamada de um $\pi$-sistema.
\end{definicao} 
 
\begin{exemplo}
	A coleção de todos os intervalos abertos da reta da 
	forma $(a,b)$ com $a<b$ é um $\pi$-sistema.
\end{exemplo}
 
    
\begin{definicao}
 \index{$\lambda$-sistema}
    Um $\lambda$-sistema é uma coleção $\mathcal{L}$ de subconjuntos
    de $\Omega$ tais que:
    \begin{itemize}
        \item[1)] $ \Omega \in \mathcal{L}$;
        \item[2)] Se $ A \in \mathcal{L}$, então $ A^c \in \mathcal{L} $;
        \item[3)] Se $ \{A_n\}$ é uma família de conjuntos de $\mathcal{L}$, dois a
        dois disjuntos, então temos que $ \cup_{n=1}^{\infty} A_n \in \mathcal{L} $.
     \end{itemize}
\end{definicao}

\begin{exercicio}
    \label{pi-lamda-sigma}
    Mostre que se $\mathcal{L}$ é um $\pi$-sistema e um $\lambda$-sistema, então
    $\mathcal{L}$ é uma $\sigma$-álgebra.   
\end{exercicio}

\begin{teorema}[Teorema $\pi -\lambda$ de Dynkin]
    \index{Teorema!$\pi-\lambda$}\label{pi-lambda}
    Se $\mathcal{L}$ é um $\lambda$-sistema contendo um $\pi$-sistema 
    $\mathcal{C}$, então $ \sigma(\mathcal{C}) \subset \mathcal{L}$.
\end{teorema}

\begin{proof}
    Seja $ \mathcal{L}(\mathcal{C}) = \cap_{i\in I} \mathcal{L}_i $, 
    onde $I$ indexa a coleção de todos os $\lambda$-sistemas 
    contendo $\mathcal{C}$. 
    Para provar o teorema, vamos demonstrar que 
    $\mathcal{L}(\mathcal{C} )$ é um $\pi$-sistema e também 
    um $\lambda$-sistema.
    Assim segue do Exercício \ref{pi-lamda-sigma} que 
    $ \mathcal{L}(\mathcal{C})$ é uma $\sigma$-álgebra.  
    Para obter a conclusão do teorema basta usar 
    o Exercício \ref{exer-sigma-gerada-contida-todas} 
    que neste caso afirma que
    $\sigma(\mathcal{C}) \subset \mathcal{L}(\mathcal{C})\subset \mathcal{L}$.
    
    Vamos mostrar primeiramente que $\mathcal{L}(\mathcal{C})$ é um um  
    $\lambda$-sistema. Como $\Omega \in \mathcal{F}$, para todo o $\lambda$-sistema  
    $\mathcal{L}_i$, segue que $\Omega \in \mathcal{L}(\mathcal{C})$. 
    Para todo $i\in I$ temos que 
    $\emptyset \in \mathcal{L}_i$, pois $\emptyset = \Omega^c \in \mathcal{L}_i$. 
    Logo $\emptyset\in \mathcal{L}(\mathcal{C})$.
    De maneira análoga para qualquer $A \in \mathcal{L}(\mathcal{C})$, temos que 
    $A \in \mathcal{L}_i$ para todo $i\in I$, de onde segue que
    $A^c \in \mathcal{L}(\mathcal{C})$.
    Considere agora uma sequência dois a dois disjunta de subconjuntos de 
    $\mathcal{L}(\mathcal{C})$, digamos $\{A_n\}$. 
    Para cada $n\ge 1$, temos que $A_n \in\mathcal{L}_i$ para todo $i\in I$ 
    o que implica que $\cup_{n=1}^{\infty} A_n \in \mathcal{L}_i$ para todo $i\in I$ 
    e portanto $\cup_{n=1}^{\infty} A_n \in  \mathcal{L}(\mathcal{C})$ e temos provado
    que $\mathcal{L}(\mathcal{C})$ é um $\lambda$-sistema. 
    
    Vamos mostrar agora que $ \mathcal{L}(\mathcal{C})$ é um $\pi$-sistema. 
    Fixado
    $A \in  \mathcal{L}(\mathcal{C})$ defina a seguinte coleção:  
    $\mathcal{L}_A := \{ B \subset \Omega : A \cap B \in  \mathcal{L}(\mathcal{C})\}$. 
    \\
    \noindent{\bf Afirmação-1}. Se $ \mathcal{L}(\mathcal{C}) \subset  \mathcal{L}_A$, para todo o 
    $A \in  \mathcal{L}(\mathcal{C})$ então $\mathcal{L}(\mathcal{C})$ é um $\pi$-sistema.

	Para ver que a Afirmação-1 é verdadeira, basta notar que fixados 
	$A,B\in\mathcal{L}(\mathcal{C})$ temos por hipótese que
	$\mathcal{L}(\mathcal{C}) \subset  \mathcal{L}_B$ e portanto 
	$A\in\mathcal{L}_B$. Segue agora da definição de $\mathcal{L}_B$ que
	$A\cap B\in \mathcal{L}(\mathcal{C})$ e portanto a afirmação está provada.
	
    
    
    Vamos mostrar agora que $ \mathcal{L}(\mathcal{C}) \subset  \mathcal{L}_A$, para todo o 
    $A \in  \mathcal{L}(\mathcal{C})$.
    Primeiro observamos que repetindo os argumentos dados acima podemos ver facilmente que 
    $\mathcal{L} _A$ é um $\lambda$-sistema. 
    Em particular, se 
    $A \in \mathcal{C}$, então $A \cap B \in \mathcal{C}$ para todo o $B \in\mathcal{C}$,
    pois $\mathcal{C}$ é fechado para a interseção finita e assim temos que 
    $\mathcal{C} \subset\mathcal{L}_A$ o que  
    implica imediatamente que $\mathcal{L}(\mathcal{C}) \subset \mathcal{L}_A $ 
    e a prova do teorema está finalmente completa.         
\end{proof}


\begin{corolario}[Unicidade]
    Sejam $P_1$, $P_2$ duas medidas de probabilidade tais que $P_1(C)=P_2(C)$ para todos os
    eventos $C$ de um $\pi$-sistema $\mathcal{C}$, então, temos que $P_1=P_2$ em
    $\sigma(\mathcal{C})$.
\end{corolario}

\begin{proof}
    Observe que $\mathcal{L}  = \{A \subset \Omega: P_1(A) = P_2(A)\}$  é um $\lambda$-sistema. 
    Por hipótese $\mathcal{L}$ contém o $\pi$-sistema $\mathcal{C}$, daí
    segue do Teorema $\pi-\lambda$ de Dynkin que $\sigma(\mathcal{C}) \subset \mathcal{L}$.
\end{proof}

\section{O Teorema da Extensão de Carathéodory}


A família de todos os subconjuntos de um conjunto $\Omega$ será 
denotada por $\mathcal{P}(\Omega)$. Esta família também é conhecida como 
conjunto das partes de $\Omega$. Observamos que $\mathcal{P}(\Omega)$ é uma 
$\sigma$-álgebra.
Em determinadas situações, quando estamos trabalhando com uma medida $\mu$, 
podemos nos livrar de várias dificuldades técnicas se esta medida esta definida
em $\mathcal{P}(\Omega)$. Acontece que na maioria das situações mais 
interessantes, isto nem sempre ocorre. Entretanto é possível definir um outro objeto 
matemático, chamado {\it medida exterior}, que está sempre definida na 
$\sigma$-álgebra das partes.
Uma medida exterior tem a vantagem de estar sempre 
definida no conjunto das partes, mas por outro lado, em geral, 
ela não é uma função $\sigma$-aditividade. 
Mas como veremos a abaixo as medidas exteriores são objetos muito 
úteis na construção de medidas sobre um determinado conjunto $\Omega$.



\begin{definicao}[Medida Exterior]
\index{Medida!exterior}
\label{definicao-medida-exterior}
	Uma função $\mu^*:\mathcal{P}(\Omega)\to [0,\infty]$ é chamada de medida exterior 
	sobre $\Omega$ se satisfaz:
	\begin{itemize}
		\item[1)] $\mu^*(\emptyset)=0$;
		\item[2)] se $A\subset B$ então $\mu^*(A)\leq \mu^*(B)$;
		\item[3)] $\mu^{*}(\cup_{i=1}^{\infty} A_n) \leq \sum_{i=1}^{\infty} \mu^{*}(A_n)$, para toda
					sequência $A_n$ ($n=1,2,\ldots$).
	\end{itemize}
\end{definicao}






A proposição seguinte é uma poderosa máquina de construir medidas exteriores.
Como o leitor pode ver as hipóteses são muito fracas. Ela nos diz que
escolhida uma coleção arbitraria $\mathcal{A}$ de subconjuntos
de um espaço $\Omega$ 
(note que não exigimos nenhuma estrutura nesta coleção nem de álgebra, $\sigma$-álgebra e etc.)
e qualquer função $\mu:\mathcal{A}\to [0,\infty]$  
tal que $\mu(\emptyset)=0$, podemos construir a partir da coleção $\mathcal{A}$ e
da função $\mu$ uma medida exterior, que será chamada de $\mu^{*}$,
definida em toda a $\sigma$-álgebra das partes de $\Omega$.
O leitor deve examinar com cuidado a definição de
$\mu^{*}$, dada abaixo, para se convencer que nem sempre $\mu^{*}$ 
é uma extensão de $\mu$. 

 
\begin{proposicao}\label{prop-med-ext}
Seja $\mathcal{A}$ uma coleção arbitrária de subconjuntos de $\Omega$ satisfazendo
apenas que $\emptyset,\Omega\in \mathcal{A}$. Seja $\mu:\mathcal{A}\to [0,\infty]$,
uma função arbitrária, satisfazendo apenas que $\mu(\emptyset)=0$.
Para todo $A\subset \Omega$, defina
\begin{equation}\label{def-medida-exterior}
\mu^{*}(A)= 
\inf \left\{
	\sum_{n}\mu(F_n): F_n\in \mathcal{A} \ \forall n,\ \ A\subset \cup_{n} F_n
\right\}.
\end{equation}
Então $\mu^{*}$ é uma medida exterior em $\Omega$.
\end{proposicao}



\begin{proof}
Já que $\emptyset\subset \emptyset \in \mathcal{A}$, temos que $\mu^{*}(\emptyset)=0$.
Se $A\subset B$, então qualquer coleção enumerável $F_n$ ($n=1,2,\ldots$) que é uma 
cobertura de $B$ (isto é, $B\subset \cup_n F_n$) é também uma cobertura de $A$.
Logo segue da definição de ínfimo que $\mu^{*}(A)\leq \mu^{*}(B)$.
Resta agora verificar que $\mu^{*}$ satisfaz a propriedade 3) de medida exterior.
Seja $A_n$ ($n=1,2,\ldots$) uma coleção enumerável de subconjuntos de $\Omega$
e $A=\cup_n A_n$. Se para algum índice $n$ temos $\mu^{*}(A_n)=\infty$,
então segue da monotonicidade $\mu^{*}$ que $\mu^{*}(A)=\infty$, e portanto 
3) é válida neste caso. Vamos assumir então que $\mu^{*}(A_n)<\infty$ 
para todo $n$. Fixe $\varepsilon>0$ arbitrário. Para cada $n$ existe
uma sequência $F_{n,k}$ ($k=1,2,\ldots$) em $\mathcal{A}$ tal que 
$A_n\subset \cup _k F_{n,k}$ e 
\begin{equation}\label{epsilon-aproximacao-med-ext}
\sum_k \mu(F_{n,k}) -\mu^{*}(A_n) < \frac{\varepsilon}{2^n},
\qquad
\text{ para todo}\ n.
\end{equation}
Certamente temos que 
$$
A\subset \bigcup_{n}\bigcup_{k} F_{n,k}.
$$
Usando a monotonicidade, em seguida, a definição de $\mu^{*}$ e por 
último a desigualdade \eqref{epsilon-aproximacao-med-ext} temos que
$$
\mu^{*}(A)
\leq 
\mu^{*}\left( \bigcup_{n}\bigcup_{k} F_{n,k} \right)
\leq
\sum_{n}\sum_{k} \mu(F_{n,k})
\leq \sum_{n}\mu^{*}(A_n) +\varepsilon.
$$
\end{proof}
 
 

\begin{definicao}[Medida $\sigma$-finita]
	Sejam $\mathcal{A}$ uma álgebra de subconjuntos de $\Omega$ e 
	$\mu:\mathcal{A}\to [0,\infty]$ uma medida.
	Se existe uma sequência $A_n$ ($n=1,2,\ldots$) tal que $\mu(A_n)<\infty$ 
	para todo $n\in\mathbb{N}$ e $\cup_{n=1}^{\infty}A_n =\Omega$, então 
	dizemos que $\mu$ é uma medida $\sigma$-finita.
\end{definicao} 
 
Precisamos apenas de mais uma definição para podermos apresentar o enunciado
preciso do Teorema da Extensão de Carathéodory.




\begin{definicao}[Conjuntos $\mu^{*}$-mensuráveis]
Dada uma medida exterior $\mu^{*}$ definida sobre as partes de um conjunto $\Omega$,
dizemos que $A\subset \Omega$ é $\mu^{*}$-mensurável se a condição abaixo é satisfeita:
\begin{equation}\label{cond-caratheodory}
\mu^{*}(E) = \mu^*(E\cap A) + \mu^*(E\cap A^c),
\qquad
\text{ para todo } \ E\subset\Omega.
\end{equation}
A condição acima é conhecida como {\it Condição de Carathéodory}\index{Condição!de Carathéodory}.
\end{definicao}




\begin{teorema}[Teorema da Extensão de Carathéodory]
\index{Teorema!da Extensão de Carathéodory}
\label{teorema-caratheodory}
Seja $\mu^*:\mathcal{P}(\Omega)\to [0,\infty]$ uma medida exterior sobre $\Omega$.
\begin{itemize}
	\item[1)] A coleção $\mathcal{M}$ de todos os conjuntos $\mu^*$-mensuráveis
				é uma $\sigma$-álgebra e a restrição de $\mu^*$ a $\mathcal{M}$
				é uma medida completa.
	\item[2)] Se $\mu^*$ é definida por \eqref{def-medida-exterior}, com $\mathcal{A}$
				sendo uma álgebra e $\mu$ uma medida em $\mathcal{A}$. Então 
				$\sigma(\mathcal{A})\subset \mathcal{M}$ e $\mu^*=\mu$ em $\mathcal{A}$.
	\item[3)] Se $\mu$ é uma medida $\sigma$-finita em uma álgebra $\mathcal{A}$ 
				então ela se estende {\bf unicamente}
				 a uma medida definida em $\sigma(\mathcal{A})$. Esta extensão 
				 é dada por $\mu^*$, definida em \eqref{def-medida-exterior}, restrita
				 a $\sigma(\mathcal{A})$.
\end{itemize}
\end{teorema}


\begin{proof}
Prova do item 1). Mostraremos primeiro que $\mathcal{M}$ é uma álgebra.
Em primeiro lugar observamos que $A=\emptyset$ satisfaz 
a condição de Carathéodory. Também é imediato verificar que se $A$ 
satisfaz a condição de Carathéodory, então $A^c$ também satisfaz esta condição. 
Pela subaditividade da medida exterior ( propriedade 3 da Definição \ref{definicao-medida-exterior}), 
para verificar a condição de Carathéodory basta mostrar que 

\begin{equation}\label{reducao-condicao-caratheodory}
\mu^*(E\cap A)+\mu^*(E\cap A^c)\leq \mu^*(E),
\qquad
\text{para todo}\ E\in\Omega.
\end{equation}
Vamos mostrar que $\mathcal{M}$ é fechado com respeito a interseções finitas. 
Para quaisquer $A,B\in\mathcal{M}$, temos
\begin{align*}
\mu^*(E) &= \mu^*(E\cap B) + \mu^*(E\cap B^c) \qquad 
\\
=& \mu^*(E\cap B\cap A)+\mu^*(E\cap B\cap A^c) + \mu^*(E\cap B^c\cap A)+ \mu^*(E\cap B^c\cap A^c)
\\
\geq & \mu^*(E\cap(B\cap A))+\mu^*(E\cap(B\cap A)^c),
\end{align*}
onde na última desigualdade, usamos que 
$(B\cap A)^c=B^c\cup A^c = (B^c\cap A) \cup (B^c\cap A^c) \cup (B\cap A^c)$
e a subaditividade de $\mu^*$. Isto mostra que a desigualdade 
\eqref{reducao-condicao-caratheodory} é satisfeita para $A\cap B$, logo 
$A\cap B\in \mathcal{M}$. Já que $\mathcal{M}$ é fechada para interseções,
podemos concluir que $\mathbb{M}$ é fechada para uniões finitas.
E portanto concluímos a prova que $\mathcal{M}$ é uma álgebra.

Próximo passo é provar que $\mathcal{M}$ é uma $\sigma$-álgebra e
que $\mu^*$ é $\sigma$-aditiva. 
Seja $B_n$ ($n=1,2,\ldots$) uma sequência de conjuntos dois a 
dois disjuntos em $\mathcal{M}$. 
Defina $C_m = \cup_{n=1}^m B_n$ ($m=1,2,\ldots$).
Vamos mostrar por indução em $m$ que 

\begin{equation} \label{caratheodory-eq1}
    \mu^*(E\cap C_m) = \sum_{n=1}^m \mu^*(E\cap B_n), 
    \qquad \text{para todo} \ E\subset\Omega.
\end{equation}

Já que $C_1=B_1$ a fórmula acima é verdadeira para $m=1$. Suponha então que 
\eqref{caratheodory-eq1} seja satisfeita para um dado $m$. Como $B_{m+1} \in  \mathcal{M}$,
para todo o $ E \subset \Omega $, temos que
\begin{align*}
    \mu^*(E \cap C_{m+1}) & = \mu^*((E \cap C_{m+1})\cap B_{m+1}) + 
    \mu^*((E \cap C_{m+1})\cap B^c_{m+1})\\
    & = \mu^*(E \cap B_{m+1}) + \mu^*(E\cap C_{m})\\
    & = \mu^*(E \cap B_{m+1}) + \sum_{n=1}^m \mu^*(E\cap B_n),
\end{align*}
onde na última igualdade usamos a hipótese de indução. 
A igualdade acima mostra que \eqref{caratheodory-eq1} 
é satisfeita para $m+1$ no lugar de $m$ e assim a indução está completa. 


Seja $A = \cup_{n=1}^{\infty}B_n $. 
Para todo $m\in\mathbb{N}$ e $E \subset \Omega$, temos que
\begin{align*}
    \mu^*(E)&= \mu^*(E\cap C_m) + \mu^*(E \cap C^c_m)
    \qquad (\text{pois  }  C_m \in \mathcal{M} ) \\
    & = \sum_{n=1}^m  \mu^*(E\cap B_n) + \mu^*(E \cap C^c_m)\\
    & \ge  \sum_{n=1}^m  \mu^*(E\cap B_n) + \mu^*(E \cap A^c),
\end{align*}
pois $ A^c \subset C^c_m$. Tomando $ m \to \infty $, ficamos com a seguinte desigualdade
\begin{equation} \label{caratheodory-eq2}
    \mu^*(E) \ge \sum_{n=1}^{\infty}  \mu^*(E\cap B_n) + \mu^*(E \cap A^c) \ge 
     \mu^*(E\cap A) + \mu^*(E \cap A^c),
\end{equation}
onde na última desigualdade usamos a propriedade 
subaditiva da medida exterior. Isso mostra que 
$ A \equiv \cup_{n=1}^{\infty} B_n \in \mathcal{M} $, 
isto é, $ \mathcal{M} $ é fechado para
uniões disjuntas contáveis. Se $A_n$, ($n=1,2,\ldots$) 
é uma sequência em $ \mathcal{M} $, 
podemos expressar $ A \equiv \cup_{n=1}^{\infty} A_n $ 
como $  A = \cup_{n=1}^{\infty}B_n$, onde
$ B_1 = A_1$ e $ B_n = A^c_1 \cap \dots \cap A^c_{n-1}\cap A_n$ para $ n \ge 2 $.
Note que a sequência $B_n$ ($n=1,2,\ldots$) definida desta forma é 
uma sequência de conjuntos dois a dois disjuntos e cada $B_n\in\mathcal{M}$. 
Então $ A \in \mathcal{M}$, provando que $ \mathcal{M}$
é fechada para uniões enumeráveis arbitrárias. 



Para provar a $\sigma$-aditividade de $ \mu^* $ em $ \mathcal{M}$, 
considere $B_n$ ($n=1,2,\ldots$) uma sequência de conjuntos 
dois a dois disjuntos em $\mathcal{M}$. 
Tomando $ E = A \equiv \cup_{n=1}^{\infty} B_n $ na primeira desigualdade em 
\eqref{caratheodory-eq2} obtemos a seguinte estimativa: 
$ \mu^*\left( \cup_{n=1}^{\infty} B_n \right)  \ge \sum_{n=1}^{\infty} \mu^*\left(  B_n \right) $. 
Usando a propriedade subaditiva de uma medida exterior concluímos que
$$ 
\mu^*\left( \cup_{n=1}^{\infty} B_n \right) = \sum_{n=1}^{\infty} \mu^*\left(  B_n \right).
$$
Assim, concluímos a prova de que $\mu^*$ é uma medida na $\sigma$-álgebra $ \mathcal{M}$. 

Vamos provar agora a última parte do item 1), isto é, 
mostrar que a medida que acabamos de obter é completa.
Sejam $N \in \mathcal{M}$ tal que $ \mu^*(N)=0$ e $A \subset N$. 
Pela monotonicidade de $\mu^*$ temos que 
$ \mu^*(E \cap A) \le \mu^*(A) \le \mu^*(N) = 0$ e $ \mu^*(E \cap A^c) \le \mu^*(E) $.
Mas estas duas desigualdades implicam
que \eqref{reducao-condicao-caratheodory} é satisfeita, provando que $ A \in \mathcal{M}$. 
O que é suficiente para concluir que a $\sigma$-álgebra $\mathcal{M} $ é $\mu^*$-completa.


Prova do item 2). Considere agora o caso em que $ \mathcal{A}$ é uma álgebra, $\mu$ é uma 
medidade em $ \mathcal{A} $, e $ \mu^* $ é a medida exterior definida em \eqref{def-medida-exterior}.
Para provar que $ \mathcal{A} \subset \mathcal{M} $, seja $ A \in \mathcal{A} $. Fixe 
$E \subset \Omega $ e $ \epsilon > 0 $ arbitrariamente. 
Invocando a definição de medida exterior \eqref{def-medida-exterior}, 
podemos afirmar que existe $ A_n \in \mathcal{A}$, 
$(n =1, 2 , \dots)$ tal que $ E \subset \cup_{n=1}^{\infty} A_n $ e 
\begin{equation}\label{estimativa-teo-extensao-carateodory-1}
  \sum_{n=1}^{\infty} \mu(A_n)-\epsilon \leq \mu^*(E).
\end{equation}
Além disso, segue da subaditividade da medida exterior e em seguida, 
da definição dada em \eqref{def-medida-exterior} as seguintes desigualdades: 
\begin{align*}
    &\mu^*(E \cap A)  \le \mu^*\left( A \cap \bigcup_{n=1}^{\infty} A_n\right) 
                      \le \sum_{n=1}^{\infty}\mu(A \cap  A_n)
	%     
     \\[0.2cm]
     \text{e} \hspace*{3cm}&
     \\[0.2cm]
     %
     &\mu^*(E \cap A^c)  \le \mu^*\left(A^c \cap \bigcup_{n=1}^{\infty} A_n\right) 
                                \le \sum_{n=1}^{\infty}\mu(A^c \cap  A_n).
\end{align*}
Somando as duas desigualdades em ambos os lados e 
usando a estimativa \eqref{estimativa-teo-extensao-carateodory-1}
obtemos 
\begin{align*}
     \mu^*(E \cap A)+ \mu*(E \cap A^c)
     \le& 
     \sum_{n=1}^{\infty}[\mu(A \cap  A_n) +  \mu(A^c \cap  A_n)]
     \\[0.2cm]
     =& 
     \sum_{n=1}^{\infty}\mu(A_n) 
     \\[0.2cm]
     \leq&
     \mu^*(E)+\varepsilon.
\end{align*}
Como $\varepsilon>0$ é arbitrário a condição \eqref{reducao-condicao-caratheodory} 
é satisfeita, provando que $ A \in \mathcal{M}$. 
Para provar que $ \mu = \mu^* $ em $ \mathcal{A} $, seja $ A \in \mathcal{A}$. Pela defininição 
\eqref{def-medida-exterior}, $ \mu^*(A) \le \mu(A) $ (tomando $ A_1 = A $ e $ A_n = \emptyset $,
para $ n \ge 2 $ por exemplo). Por outro lado, $ \mu(A) \le \sum_{n=1}^{\infty} \mu(A_n)  $ para toda
a sequência $ A_n \in \mathcal{A} $ ($  n \ge 1$) tal que $ A \subset \cup_{n=1}^{\infty} A_n $, então
temos que  $ \mu^*(A) \ge \mu(A) $ (por subaditividade de $\mu$ em $ \mathcal{A} $).
Mostrando que $ \mu(A) = \mu^*(A)$.

Prova do item 3). 
Suponha que $ \mu$ seja $\sigma$-finita 
na álgebra $\mathcal{A}$ e seja
$\mu^*$ sua extensão a $\sigma(\mathcal{A}) \subset \mathcal{M}$, 
dada pelo item 2) deste teorema. Já que $\mu$ é $\sigma$-finita, podemos
encontrar uma sequência $A_n$ ($n=1,2,\ldots$) de conjuntos
dois a dois disjuntos tal que 
$A_n\in \mathcal{A}$, $\mu(A_n)<\infty$, para todo $n\in\mathbb{N}$ e 
$\Omega = \cup_{n=1}^{\infty} A_n$. 
Seja $\nu$ uma extensão da medida $\mu$, definida em $\sigma(\mathcal{A})$.
Fixado $n\in\mathbb{N}$, vamos mostrar primeiro que $\mu^{*} = \nu$, 
em todos os conjuntos da coleção 
$
A_n\cap \sigma(\mathcal{A}) \equiv \{ A_n\cap A: A\in\sigma(\mathcal{A})\}.
$
Para provar este fato vamos mostrar que a coleção 
$
\mathcal{C} = \{ A\in \sigma(\mathcal{A}) : \nu(A_n\cap A) = \mu^*(A_n\cap A)\}
$
é um $\lambda$-sistema contendo $\mathcal{A}$ (que é um $\pi$-sistema). 
De fato, se $A \in\mathcal{C}$ então temos que: 
\begin{align*}
	&\nu(A_n) =
	\nu((A_n\cap A^c) \cup (A_n\cap A)) 
	= 
	\nu(A_n\cap A^c)+ \nu(A_n\cap A)
\\
\text{e}\hspace*{1cm}&
\\
	&\mu^*(A_n) =
	\mu^*((A_n\cap A^c) \cup (A_n\cap A)) 
	= 
	\mu^*(A_n\cap A^c)+ \mu^*(A_n\cap A).
\end{align*}
Observando que o lado esquerdo de ambas igualdades acima são iguais
($\nu$ e $\mu^*$ são extensões de $\mu$ em $\mathcal{A}$) 
e que as segundas parcelas do lado direito também são iguais 
já que $A\in\mathcal{C}$, concluímos que $\mu^*((A_n\cap A^c)= \nu(A_n\cap A^c)$ 
e portanto $A^c\in\mathcal{C}$. 
Como o $\emptyset\in\mathcal{C}$ segue do fato que acabamos de provar que
$\Omega\in\mathcal{C}$. 
Seja $B_m$ ($m=1,2,\ldots$) uma sequência de conjuntos em 
$\mathcal{C}$ dois a dois disjuntos. Segue da $\sigma$-aditividade
de $\mu^*$ e $\nu$ que 
$$
\nu\left( A_n \cap \big(\cup_{m=1}^{\infty}B_m\big) \right)
=
\sum_{m=1}^{\infty} \nu\left( A_n \cap B_m \right)
=
\sum_{m=1}^{\infty} \mu^*\left( A_n \cap B_m \right)
=
\mu^*\left( A_n \cap \big(\cup_{m=1}^{\infty}B_m\big) \right).
$$
O que encerra a prova de que $\mathcal{C}$ é um $\lambda$-sistema.
Como $\mathcal{C}$ é um $\lambda$-sistema que contém o $\pi$-sistema
$\mathcal{A}$ segue do Teorema $\pi-\lambda$ de Dynkin que 
$\sigma(A)\subset \mathcal{C}$.


Para finalizar a prova, basta notar que dado $A\in \sigma(\mathcal{A})$, 
podemos escrever
$A=\cup_{n=1}^{\infty} (A\cap A_n)$.
Agora, usando a continuidade das medidas $\nu$ e $\mu^*$,
temos que  
\begin{align*}
\nu(A) = \nu\left( \cup_{n=1}^{\infty} (A\cap A_n) \right)
	   = \lim_{n\to\infty} \nu(A\cap A_n)
	   =& \lim_{n\to\infty} \mu^{*}(A\cap A_n)
	   \\[0.2cm]
	   =& \mu^*( \cup_{n=1}^{\infty} (A\cap A_n) )
	   \\[0.2cm]
	   =& \mu^*(A).
\end{align*}
\end{proof}
	
% Aula 3
\chapter[Aula 3]{Medida de Lebesgue em $\mathbb{R}$}
\chaptermark{}






\section*{Medida de Lebesgue}


Queremos construir em $\mathbb{R}$ uma medida que 
generalize a noção de comprimento de  
um intervalo, isto é, queremos uma medida  
$\mu$ definida em alguma classe razoável de 
subconjuntos de $\mathbb{R}$ tal que se I é um intervalo
com fecho $\overline{I}=[a,b]$ então $\mu(I)=b-a$ 
e se $A=\bigsqcup_{j=1}^nI_j$ onde cada $I_j$ 
é um intervalo do tipo $(a_j,b_j), [a_j,b_j],(a_j,b_j]$ ou $[a_j,b_j)$,  
então $\mu(A)=\sum_{j=1}^n(b_j-a_j)$ e
além do mais $\mu(\emptyset)=0$.



Pois bem, dado $A\subset \mathbb{R}$ 
definimos a medida exterior de $A$ e 
denotamos $\tn{Leb}^{*}(A)$ como sendo 
$$
\tn{Leb}^{*}(A)=\inf 
				\left\{
					\sum_{j=1}^\infty \mu(I_j):~ A\subset \bigcup_{j=1}^\infty I_j 
				\right\}
$$
onde o ínfimo é tomado sobre todas as coberturas de $A$ por intervalos fechados.  

Segue da Proposição \ref{prop-med-ext} que 
$\tn{Leb}^{*}:\mathcal{P}(\mathbb{R})\to[0,\infty]$ é de fato uma medida exterior.

\begin{lema}\label{LEB 2}
$\tn{Leb}^{*}([a,b])=b-a$
\end{lema}

\begin{proof}
É claro que $\tn{Leb}^{*}([a,b])\leq b-a$, 
precisamos apenas provar a desigualdade oposta. 
Comecemos tomando uma cobertura de 
$[a,b]\subset \bigcup_{k\in \mathbb{N}} (c_k,d_k)$  
por intervalos abertos, por compacidade podemos considerar 
apenas uma subcobertura finita de $[a,b]$, 
$$
[a,b]\subset \bigcup_{k=1}^M(c_k,d_k)
$$
suponha sem perda de generalidade que esta subcobertura seja minimal, 
isto é, se tirarmos qualquer intervalo que compõe a subcobertura 
não teremos mais uma subcobertura. Para facilitar o raciocínio 
ordenaremos os intervalos de modo que 
$$
c_0<c_1<\cdots<c_{M-1}
$$
note que a minimalidade da subcobertura garante então que 
$$
d_0<d_1<\cdots<d_{M-1}
$$
para fixar raciocínio  suponha que tivessemos $d_2<d_1$ 
então teríamos que $(c_2,d_2)\subset (c_1,d_1)$ o 
que contraria a minimalidade da cobertura. 
({\red um desenho aqui deixa tudo óbvio}) 
Agora defina 
$e_0=a$, $e_1=\frac{c_1+d_0}{2},$ $e_2=\frac{c_2+d_1}{2}$,
$e_{M-1}=\frac{c_{M-1}+d_{M-2}}{2}$ e $e_{M}=b.$  
Agora note que  para cada $k$, 
tem-se $[e_k,e_{k+1}]\subsetneqq [c_k,d_k]$ assim 
$$
\sum_{k=0}^{M-1} (d_k-c_k)>\sum_{k=0}^{M-1} (e_{k+1}-e_k)=b-a
$$ 
donde $\tn{Leb}^{*}([a,b])\geq b-a$.

\end{proof}

\begin{lema}

$\tn{Leb}^{*}((a,b))=b-a$
\end{lema}


\begin{proof}
Como o intervalo $(a,b)$ está contido no seu 
fecho $[a,b]$ portanto temos que  $\tn{Leb}^{*}((a,b))\leq b-a.$ 
Agora considere um intervalo $[a+\epsilon,b-\epsilon]$,
para algum $\epsilon>0$, 
contido em $(a,b)$ note que qualquer cobertura de 
$(a,b)$ por intervalos fechados é também uma cobertura 
de $[a+\epsilon,b-\epsilon]$ de modo que 
$$
b-a-2\epsilon= \tn{Leb}^{*}([a_0,b_0])\leq  \tn{Leb}^{*}((a,b))\leq b-a
$$
como $\epsilon$ é arbitrário segue que  $\tn{Leb}^{*}((a,b))=b-a$.
\end{proof}
Usando a mesma técnica da demonstração acima pode-se mostrar 
resultados análogos considerando intervalos  
da forma $[a,b)$ e $(a,b]$.





\begin{lema}\label{LEB 3}
Seja $A\subset \mathbb{R}$,  aberto e $A=\bigsqcup_{k=1}^M(a_k,b_k)$ então temos 
$$
\tn{Leb}^{*}(A)=\sum_{k=1}^M(b_k-a_k).
$$
\end{lema}

\begin{proof}
Segue diretamente da defininção que 
$\tn{Leb}^{*}(A)\leq\sum_{k=1}^M(b_k-a_k)$. 
Suponha que vale a desigualdade estrita, 
isto é, $\tn{Leb}^{*}(A)<\sum_{k=1}^M(b_k-a_k)$. 
Seja $a=\inf A$ e $b=\sup A$
e considere o conjunto $[a,b]\setminus A=\bigsqcup_{k=1}^{N}[c_l,d_l]$ 
então temos
$$
[a,b]=\left(~\bigsqcup_{k=1}^M(a_k,b_k)~\right)
\cup \left(~\bigsqcup_{l=1}^{N}(c_l,d_l)~\right)
$$
então pela propriedade da medida exterior 
\begin{eqnarray*}
\tn{Leb}^{*}([a,b])
&\leq &
\tn{Leb}^*\left(~\bigsqcup_{k=1}^M(a_k,b_k)~\right)
+\tn{Leb}^*\left(~\bigsqcup_{l=1}^{N}(c_l,d_l)~\right)
\\
&
<
&
\sum_{k=1}^{M} (b_k-a_k)+\sum_{l=1}^N(d_l-c_l)
\\
&
=
&
b-a
\end{eqnarray*}
contradizendo o lema (\ref{LEB 2}).

\end{proof}



\begin{proposicao}
Todo conjunto aberto $A\subset \mathbb{R}$ pode ser 
escrito, de modo único, como uma união enumerável de 
intervalos dois a dois disjuntos.
\end{proposicao}


\begin{proof}
Seja $x\in A$  e considere os números $a_x$ e $b_x$ 
definidos como a seguir
$$
a_x
=
\inf\{x;~x<a~\tn{e}~(x,a)\subset A\},~~b_x
=
\sup \{x;~x>a~\tn{e}~(a,x)\subset A\}
$$
note que  tanto $a_x$ quanto $b_x$ estão bem definidos 
pois como $A$ é  aberto sempre existe um intervalo $I$ 
satisfazendo $a\in I\subset A$. 
Então para cada $x\in A$ considere o intervalo 
$I_x=(a_x,b_x)$. Sejam $x,y\in A$, afirmamos que ou $I_x=I_y$ 
ou $I_x\cap I_y=\varnothing.$ 
De fato, suponha que $I_x\cap I_y\neq \varnothing$ 
então a reunião $I_x\cup I_y$ é um intervalo que 
contém tanto $x$ quanto $y$ de modo que temos, 
por definição que $  I_ x\cup I_y\subset I_y$ 
e $I_x\cup I_y\subset I_x$ portanto 
$$
I_x\subset I_x\cup I_y\subset I_y~~ I_y\subset I_x\cup I_y\subset I_x
$$
donde $I_x=I_y$.



Para ver que a família $\{I_x\}_{x\in A}$ é 
contável escolha em cada intervalo $I_x$ um 
racional $r_x$ e considere a aplicação  
$f:\{I_x,x\in A\}\to \mathbb{Q}$ definida por 
$f(I_x)=r_x$ e note que $f$ é injetiva, com efeito se 
$r_y=f(I_y)=f(I_x)=r_x$ temos que 
$I_x\cap I_y\neq \varnothing$ donde $I_x=I_y$ e portanto $r_x=r_y$.

\end{proof}




\begin{lema}\label{LEB 4}
Seja $A\subset \mathbb{R}$,  aberto e 
$A=\bigsqcup_{k\in \mathbb{N}}(a_k,b_k)$ então temos 
$$
\tn{Leb}^{*}(A)=\sum_{k\in \mathbb{N}}(b_k-a_k).
$$
\end{lema}

\begin{proof}
Temos pela proposição  que 
$\tn{Leb}^{*}(A)\leq \sum_{k\in \mathbb{N}}(b_k-a_k).$ 
Suponha que tenhamos na verdade 
$\tn{Leb}^{*}(A)< \sum_{k\in \mathbb{N}}(b_k-a_k)$ 
segue dai que existe $K\in \mathbb{N}$ tal que 
$$
\tn{Leb}^{*}(A)< \sum_{k<K}(b_k-a_k)
$$
 por outro lado pelo lema (\ref{LEB 3}) 
 que $\sum_{k<K}(b_k-a_k)=\tn{Leb}^{*}(A_0)$ 
 onde $A_0=\bigcup_{k<K}(a_k,b_k)$.
Temos então uma contradição pois 
$A_0\subset A$ e $\tn{Leb}^{*}(A)<\tn{Leb}^{*}(A_0).$
\end{proof}

Note que em última análise o resultado acima 
vale para outras uniões enumeráveis disjuntas 
de intervalos não necessariamente abertos, 
isto é, vale para uniões disjuntas 
$A=\bigsqcup_{n\in\mathbb{N}} I_n$ onde 
cada intervalo $I_n$ pode ser qualquer um 
dos intervalos $[a,b),[b,a)$ ou $[a,b]$. 



%Veja que estamos considerando a semi-álgebra $\mathcal{S}$ dos cubos fechados de $\mathbb{R}^d$ e a medida finitamente aditiva $\mu$ caracterizada por 
%$\mu(C)=|C|$ onde $C$ é um cubo fechado.  Então $\tn{Leb}^{*}$ é a medida exterior induzida por esta medida finitamente aditiva.

\begin{proposicao}
Sejam $A_1, \ldots, A_n, \ldots$ elementos da semi-álgebra 
$\mathcal{S}$ tais que $A=\bigcup_{i=1}^{\infty}A_n$


\end{proposicao}




\begin{proposicao}
Seja $A\subset \mathbb{R}$ então 
$ 
\tn{Leb}^{*}(A)
=
\inf_{\mathcal{O}\supset A}\tn{Leb}^{*}(\mathcal{O})
$ 
onde $\mathcal{O}$ é um aberto contendo $A$.
\end{proposicao}


\begin{proposicao}
Sejam $A,B\subset \mathbb{R}$ com $\tn{dist}(A,B)>0$ 
então 
$
\tn{Leb}^{*}(A\cup B)
=
\tn{Leb}^{*}(A)+\tn{Leb}^{*}(B)
$.
\end{proposicao}



\subsubsection{Uma construç\~ao autom\'atica da medida de Lebesgue}


Considere álgebra $\mathcal{A}$ consistindo dos 
intervalos e uniões de intervalos de   $\mathbb{R}.$ 
Seja  $\mu$ a  função de cojunto  nessa álgebra 
caracterizada do seguinte modo, se $[a,b]$ é um 
intervalo  então $\mu([a,b])=b-a$ 
e se $A=\cup_{i=1}^n(a_i,b_i)$ então 
$\mu(A)=\sum_{i=1}^n(b_i-a_i).$  
Note que $\mu$ é claramente uma medida 
finitamente aditiva em $\mathcal{A}$, 
além do mais o Lema (\ref{LEB 4}) mostra 
que $\mu$ é $\sigma$-aditiva. 
O teorema de extensão de Caratheodory diz que a 
medida $\mu$ se estende a uma medida definida na 
$\sigma$-álgebra  $\mathcal{M}$  
caracterizada pela condição de Caratheodory. 
Além do mais se considerarmos a restrição 
da medida obtida  à $\sigma$-álgebra
$\sigma(\mathcal{S})\subset \mathcal{M}$ 
esta restrição é única. A medida que é a 
extensão de $\mu$ é dita a medida de Lebesgue 
de $\mathbb{R}$ e denotada por $\tn{Leb}$, 
a sigma álgebra $\mathcal{M}$ nesse caso particular
é dita a $\sigma$-álgebra de Lebesgue e é denotada por $\mathcal{L}$.






































\subsection*{Uma outra construç\~ao da medida de Lebesgue}













\subsubsection{Caratheodory mensur\'avel e Lebesgue mensur\'avel}




A condição que define mensurabilidade no sentido de 
Caratheodory parece muito misteriosa, ou em 
um portugês mais claro 
parece ``tirada do chapeu", o que está por 
trás disso são anos de pesquisa árdua. 
Entretanto em $\mathbb{R}^d$ é possível definir 
mensurabilidade de maneira bastante razoável. 
Diremos que $A\subset \mathbb{R}^d$ é Lebesgue  
mensurável quando para cada  $\epsilon>0$ 
existe um aberto $\mathcal{O}\supset A$ 
tal que $\tn{Leb}^{*}(\mathcal{O}\setminus A)\leq \epsilon$. 
Esta definição parece portanto bastante razoável, 
um conjunto $A\subset \mathbb{R}^d$ é mensurável 
se sua medida exterior  pode ser bem aproximada 
pela medida exterior de abertos contendo $A$.

A grande pergunta então é se o conceito de 
mensurabilidade dado acima coincide com o 
conceito de mensurabilidade 
no sentido dado por Caratheodory. 
Bom  a resposta é sim como mostraremos.



Seja $A\subset \mathbb{R}^d$ e suponha que $A$ 
seja Lebesgue mensurável. 
Note que todo aberto 
é Caratheodory mensurável. 
Dado $\epsilon>0$ seja $\mathcal{O}$ um 
conjunto aberto contendo $A$ 
tal que 
$\tn{Leb}^{*}(\mathcal{O}\setminus A)\leq \epsilon$  
então dado $E\subset \mathbb{R}^d$ 
arbitrário temos 
\begin{eqnarray*}
\tn{Leb}^{*}(E\cap A)+\tn{Leb}^{*}(E\cap A^c)
&\leq& 
\tn{Leb}^{*}(E\cap \mathcal{O})+\tn{Leb}^{*}(E\cap A^c)
\\
&
\leq 
&
\tn{Leb}^{*}(E\cap\mathcal{O})+
\tn{Leb}^{*}(E\cap (A^c\setminus \mathcal{O}^c))
\\
&
+
&
 \tn{Leb}^{*}(E\cap \mathcal{O}^c)
 \\
 &
 \leq 
 &
\tn{Leb}^{*}(E)+\epsilon
\end{eqnarray*}
onde usamos acima que 
$
A\cap E\subset \mathcal{O}\cap E,$ $A^c
=
(A^c\setminus \mathcal{O}^c)\cup \mathcal{O}^c
$ 
e que 
$
A^c\setminus \mathcal{O}^c=\mathcal{O}\setminus A$ como 
$\epsilon>0$ é arbitrário concluímos que
$$
\tn{Leb}^{*}(E)\geq \tn{Leb}^{*}(E\cap A)+\tn{Leb}^{*}(E\cap A^c)
$$
donde $A$ é Caratheodory mensurável.


A recíproca é um tanto mais simples. De fato seja 
$A\subset \mathbb{R}^d$ Caratheodory mensurável 
tal que $\tn{Leb}^{*}(A)<\infty$ dado $\epsilon>0$  
existe uma família $\{C_i\}_{i=1}^{\infty}$ de 
cubos fechados tais que 
$A\subset \bigcup_{i=1}^{\infty}C_i$ e 
$\sum_{i=1}^{\infty}|C_i|\leq \tn{Leb}^{*}(A)+\epsilon$.
Agora para cada $i$ consideremos um cubo $\tilde{C}_i$ aberto 
tal que $|\tilde{C}_i|<|C_i|+\epsilon/2^{i}.$
Como cada $\tilde{C}_i$ é Caratheodory mensurável sua 
união enumerável também o é de modo que podemos escrever 
\begin{eqnarray*}
\tn{Leb}^{*}(\bigcup_{i=1}^{\infty} \tilde{C}_i\setminus A)
&=&
\tn{Leb}^{*}(\bigcup_{i=1}^{\infty} \tilde{C}_i)-\tn{Leb}^{*}(A)
\\
&
=
&
\sum_{i=1}^{\infty}|\tilde{C}_i|-\tn{Leb}^{*}(A)
\\
&
\leq 
&
\sum_{i=1}^{\infty}|C_i|-\tn{Leb}^{*}(A)+\sum_{i=1}^{\infty}\epsilon/2^i
\\
&
\leq 
&
2\epsilon
\end{eqnarray*}
então tomando  $\mathcal{O}=\bigcup_{i=1}^{\infty}C_i$ 
temos que $\tn{Leb}^{*}(\mathcal{O}\setminus A)\leq \epsilon$.
 
Fazendo um apanhado da pequena discussão acima, 
concluímos que a noção de conjunto mensurável no 
sentido de Caratheodory quando partimos da semi-álgebra 
dos retângulos em $\mathbb{R}^d$ e a noção de mensurabilidade
no sentido de Lebesgue são as mesmas, 
em particular provamos que os mensuráveis no 
sentido de Lebesgue também formam uma sigma álgebra.




\subsubsection{Uma construção artesanal da medida de Lebesgue}

A seguir refaremos em $\mathbb{R}^d$ todo o processo 
de extensão da medida finitamente aditiva feita na 
secção anterior para a única diferença e que 
usaremos apenas o conceito de mensurabilidade 
no sentido de Lebesgue.
É claro que do ponto de vista 
rigoroso toda a construção que virá é 
desnecessária tendo em vista o que já foi discutido acima, 
por outro lado do ponto de vista didático esta 
discussão será  muito útil.









\begin{definicao}
Um subconjunto $A\subset \mathbb{R}^d$ é dito 
Lebesgue mensurável, ou simplesmente mensurável,  
se para cada $\epsilon>0$ existe um 
conjunto aberto $\mathcal{O}\subset A$ 
tal que $\tn{Leb}^{*}(\mathcal{O}\setminus A)\leq \epsilon$.
\end{definicao}


\begin{proposicao}
Todo conjunto aberto $A\subset \mathbb{R}^d$ é mensurável.
\end{proposicao}






\begin{proposicao}
Se $\tn{Leb}^{*}(A)=0$ então $A$ é  mensurável, em particular 
todo subconjunto de $A$ é mensurável.
\end{proposicao}
\begin{proof}
Temos pela proposição (??) que dado $\epsilon>0$ existe 
um aberto $\mathcal{O}\supset A$ tal que 
$\tn{Leb}^{*}(\mathcal{O})\leq \epsilon$, 
como $\mathcal{O}\setminus E\subset \mathcal{O}$ 
temos por monotonicidade da medida exterior 
que $\tn{Leb}^{*}(\mathcal{O}\setminus E)\leq \epsilon$. 
\end{proof}





\begin{proposicao}
Toda  união enumerável de conjuntos mensuráveis é mensurável.
\end{proposicao}

\begin{proof}
Seja $A=\bigcup_{j=1}^{\infty} A_j$ onde cada $A_j$ é mensurável. 
Dado $\epsilon>$ temos para cada $j$ que existe um 
aberto $\mathcal{O}_j\supset A_j$ e 
tal que $\tn{Leb}^{*}(\mathcal{O}_j\setminus A_j)\leq \epsilon/2^j$. 
Temos então que  a união 
$\mathcal{O}=\bigcup_{j=1}^{\infty}\mathcal{O}_j$ 
é um conjunto aberto,   $A\subset \mathcal{O}$ e
que 
$$
\mathcal{O}\setminus A=\bigcup_{j=1}^{\infty} \mathcal{O}_j\setminus A_j
$$
por monotonicidade e subaditividade que 
$$
\tn{Leb}^{*}(\mathcal{O}\setminus A)
\leq 
\sum_{n=1}^\infty \tn{Leb}^{*}(\mathcal{O}_j\setminus A_j)
\leq 
\sum_{n=1}^\infty \epsilon/2^j=\epsilon.
$$

\end{proof}


\begin{proposicao}
Conjuntos fechados de $\mathbb{R}^d$ são mensuráveis.
\end{proposicao}





\begin{proposicao}
Complementos de conjuntos mensuráveis são mensuráveis
\end{proposicao}


\begin{proposicao}

\end{proposicao}
  



% Aula 4
\chapter[Aula 4]{Espaços de Probabilidade}
\chaptermark{}





\section{Definições e Propriedades Básicas}

Um espaço de probabilidade é um tripla 
$(\Omega,\F,\P)$ onde 
\begin{itemize}
	\item 
	$\Omega$ é um conjunto chamado de espaço amostral.

	\item
	$\F$ uma $\sigma$-álgebra de subconjuntos de $\Omega$.
	Os elementos de $\F$ são chamados de eventos.
	
	\item $\P$ é uma medida de probabilidade, isto é, 
	uma função com domínio em $\F$ tomando valores em 
	$[0,1]$ satisfazendo 
		\begin{itemize}
			\item[i)] 
			$\P(E)\geq 0$ para todo $E\in F$.

			\item[ii)] 
			$P$ é $\sigma$-aditiva: Para qualquer
			sequência $\{E_n\}$ mutuamente disjunta 
			de eventos em $\F$ temos 
				\[
					\P\left( \bigcup_{n=1}^{\infty} E_n \right)
					=
					\sum \P(E_n).
				\]
			\item[iii)] 
			$\P(\Omega)=1$.
		\end{itemize}
\end{itemize}



Seguem das propriedade $\textrm{ii})$ e $\textrm{iii})$ acima que 
para qualquer evento $E\in\F$ vale a seguinte identidade
	\[
		P(E^c)=1-P(E).
	\]
De fato, 
	\[
		1=P(\Omega)= P(E\cup E^c)=P(E)+P(E^c).
	\]
Como toda medida de probabilidade é em particular uma medida então 
temos que 
	\[
		\P(\emptyset)=0.
	\] 
\begin{proposicao} Seja $(\Omega,\F,\P)$ um espaço de probabilidade.
Se $A,B$ são dois eventos arbitrários então 
	\[
		\P(A\cup B) = \P(A)+\P(B)-\P(A\cap B).
	\]
\end{proposicao}
\begin{proof}
Primeiro observamos que 
%	
	\begin{align*}
		\P(A) = \P(A\cap B^c)+\P(A\cap B) 
		\\
		\P(B) = \P(B\cap A^c)+\P(B\cap A) 
	\end{align*}
Note que podemos escrever $A\cup B = (A\cap B^c) \cup (B\cap A^c) \cup (A\cap B)$.
Tomando probabilidade em ambos lados desta igualdade 
e usando as identidades acima, temos
	\begin{align*}
	\P(A\cup B) 
	&=
	\P(A\cap B^c) + \P(B\cap A^c) + \P(A\cap B)
	\\
	&=
	[\P(A)-\P(A\cap B)] + [\P(B)-\P(A\cap B)] + P(A\cap B)
	\\
	&=
	\P(A)+\P(B)-\P(A\cap B).
	\end{align*}


\end{proof}

\begin{teorema}[Fórmula de Inclusão-Exclusão]
Para quaisquer eventos $E_1,\ldots,E_n$ temos 
	\begin{align*}
		\P\left( \bigcup_{j=1}^{n} E_j \right)
		&=
		\sum_{j=1}^n \P(E_j)
		-
		\sum_{1\leq i<j\leq n} \P(E_i\cap E_j) 
		+
		\sum_{1\leq i<j<k\leq n} \P(E_i\cap E_j\cap E_k)-\ldots
		\\[0.3cm]
		&\hspace*{0.5cm} \ldots(-1)^n \P(E_1\cap\ldots\cap E_n).
	\end{align*}
\end{teorema}

\begin{proof}
 Podemos provar o teorema por indução. Para $n=2$ o teorema
 segue diretamente da proposição acima. Para dar o passo de 
 indução basta escrever $\cup_{j=1}^{k+1} E_j$ 
 como $\cup_{j=1}^{k}E_j\cup E_{k+1}$, usar o caso $n=2$ 
 para calcular a probabilidade desta decomposição, em seguida usar 
 a hipótese de indução para 
 \[	
 \P \left( \bigcup_{j=1}^{k} E_j \right)
 \qquad
 \text{e}
 \qquad
 \P \left( \bigcup_{j=1}^{k} (E_j\cap E_{k+1}) \right)
 \]
 e reorganizar os termos de ambas as somas. Observe 
 por exemplo, que as probabilidades de interseções de pares 
 com $E_j$'s com índice menores ou iguais a $k$ virão 
 da primeira probabilidade e interseção de pares de
 eventos $E_j$'s onde um dos eventos é $E_{k+1}$ virá 
 da segunda soma e assim por diante.
\end{proof}

Uma aplicação deste Teorema é a obtenção das
chamadas {\bf Desigualdades de Bonferroni} \index{Desigualdade!Bonferroni}.
A prova destas desigualdades segue simplesmente de negligenciarmos 
os resto. Dois exemplos de tais desigualdades são mostrados
abaixo:
\[ 
\begin{array}{c}
	\displaystyle
	\P \left( \bigcup_{j=1}^{n} E_j \right) 
	\leq 
	\sum_{j=1}^n P(E_j)
	%
	%
	\\[0.5cm]
	%
	%
	\displaystyle
	\P \left( \bigcup_{j=1}^{n} E_j \right) 
	\geq
	\sum_{j=1}^n \P(E_j) - \sum_{1\leq i<j\leq n} \P(E_i\cap E_j)	
\end{array}
\]

Outra propriedade importante é a {\bf Monotonicidade} de $\P$, isto é, 
para quaisquer eventos $A,B$ tais que $A\subset B$ temos 
$\P(A)\leq \P(B)$. De fato, 
$\P(B)=\P(B\cap A)+\P(B\cap A^c)= \P(B)+\P(B\setminus A)\geq \P(A)$.

Toda medida de probabilidade é $\sigma$-subaditiva, ou seja, 
para quaisquer eventos $E_1,E_2,\ldots$ temos 
	\[
		\P\left( \bigcup_{n=1}^{\infty} E_n\right)
		\leq 
		\sum_{n=1}^{\infty} \P(E_n).
	\]
Para verificar que esta propriedade é verdadeira basta observar que 
\[
	\bigcup_{n=1}^{\infty} E_n = E_1+E_1^c\cap E_2 + E_1^c\cap E_2^c\cap E_3 +\ldots
\]
aplicar a $\sigma$-aditividade de $\P$ e em seguida usar a monotonicidade 
obtendo
	\begin{align*}
	\P\left( \bigcup_{n=1}^{\infty} E_n\right)
	&=
	\P(E_1)+\P(E_1^c\cap E_2) + \P(E_1^c\cap E_2^c\cap E_3) +\ldots
	\\
	&\leq
	\P(E_1) + \P(E_2) + \P(E_3) +\ldots
	\end{align*}
	
\begin{teorema}[Continuidade da Probabilidade]
	Sejam $(\Omega,\F,\P)$ um espaço de Probabilidade e
	$\{E_n\}$ uma sequência de eventos neste espaço.
	\begin{enumerate}
		\item
		Se $E_n\nearrow E$ então 
		$\displaystyle\lim_{n\to\infty}\P(E_n)=\P(E)$.
		\item 
		Se $E_n\searrow E$ então 
		$\displaystyle\lim_{n\to\infty}\P(E_n)=\P(E)$.		
	\end{enumerate}
\end{teorema}

\begin{proof}
Vamos provar o primeiro item do Teorema. 
Suponha que $E_1\subset E_2\subset \ldots\subset E_n\subset \ldots$
e defina a seguinte sequência de eventos 
\[
B_1= E_1,\ B_2= E_2\setminus E_2,\ \ldots,\  B_n=E_n\setminus E_{n-1}.
\]
Por construção a sequência $\{B_n\}$ é uma sequência de eventos
mutuamente disjunta e além do mais 
\[
\bigcup_{j=1}^n B_j = E_n, \quad 
\bigcup_{j=1}^{\infty} B_j = \bigcup_{j=1}^{\infty} E_j = E.
\]
Da $\sigma$-aditividade segue que 
\begin{align*}
	\P(E)
	&=
		\P\left( \bigcup_{n=1}^{\infty} E_n\right)
		=
		\sum_{n=1}^{\infty}\P(B_n)
		=
		\lim_{n\to\infty} \sum_{j=1}^n \P(B_j)
	\\[0.3cm]
	&=
		\lim_{n\to\infty}\P\left( \bigcup_{j=1}^{n} B_j\right)
		=
		\lim_{n\to\infty}\P\left( E_n\right).
\end{align*}

Para provar o segundo item note que se $E_n\searrow E$ 
então $E_n^c \nearrow E^c$. Aplicando o resultado que 
acabamos de provar temos que 
	\[
		1-\P(E)
		=
		\P(E^c)		
		=
		\lim_{n\to\infty} \P(E_n^c)
		=
		\lim_{n\to\infty} [1-\P(E_n)]		
	\]
de onde segue o resultado.
\end{proof}

Vamos provar agora um lema de continuidade para sequências 
que não são necessariamente monótonas. 
Este lema será chamado de lema de Fatou para funções indicadoras
e a razão para isto ficará clara futuramente.

\begin{lema}[Lema de Fatou para Funções Indicadoras]
Seja $(\Omega,\F,\P)$ um espaço de probabilidade e 
$\{E_n\}$ uma sequência de eventos.
\begin{enumerate}
	\item
	$\displaystyle
	\P(\liminf_{n\to\infty} E_n)
	\leq 
	\liminf_{n\to\infty} \P(E_n)
	\leq
	\limsup_{n\to\infty} \P(E_n)
	\leq
	\P(\limsup_{n\to\infty} E_n)$.
	
	\item Se $E_n\to E$ então $\P(E_n)\to \P(E)$.
\end{enumerate}
\end{lema}


\begin{proof}
Primeiro vamos assumir que 1 é válido e vamos mostrar 2.
Suponha que $E_n\to E$. Então 
\[
 \liminf_{n\to\infty} E_n = \limsup_{n\to\infty} E_n = E.
\]
Aplicando 1 temos que 
	\begin{align*}
	\P(E)
	&=\displaystyle
		\P(\liminf_{n\to\infty} E_n)
		\leq 
		\liminf_{n\to\infty} \P( E_n)
		\leq
		\limsup_{n\to\infty}\P(E_n)
		\leq
		\P(\limsup_{n\to\infty} E_n)
		=
		P(E).
	\end{align*}
O que prova que $\P(E_n)\to P(E)$.

Passamos agora a prova do item 1. Por definição 
do $\liminf$, continuidade da probabilidade
temos 
\begin{align*}
	\P(\liminf_{n\to\infty} E_n)
	&=
	\P\left( \lim_{n\to\infty} \bigcap_{k\geq n}E_k \right)
	\\
	&=
	\lim_{n\to\infty} \P\left( \bigcap_{k\geq n}E_k \right)
\end{align*}
Por monotonicidade temos que 
	\[
		\P\left( \bigcap_{k\geq n}E_k \right)
		\leq 
		\P(E_n).
	\] 
Assim se tomamos $\liminf$ em ambos os lados e usamos as igualdades acima,
obtemos a seguinte desigualdade
	\begin{align*}
	\P(\liminf_{n\to\infty} E_n)
	\leq 
	\liminf_{n\to\infty}\P(E_n).
\end{align*} 

De maneira análoga, temos 
\begin{align*}
	\P(\limsup_{n\to\infty} E_n)
	&=
	\P\left( \lim_{n\to\infty} \bigcup_{k\geq n}E_k \right)
	\\[0.2cm]
	&=
	\lim_{n\to\infty} \P\left( \bigcup_{k\geq n}E_k \right)
	\\[0.3cm]
	&\geq 
	\limsup_{n\to\infty} \P(E_n).
\end{align*}
\end{proof}


%%%%%%%%%%%%%%%%%% Referências e Índice Remissivo
\backmatter
%
%\chapter{Referências}
%\include{./referencias/referencias}

%\printindex
%\chapter{Índice de Símbolos}


\end{document}
